\documentclass[11pt,a4paper]{article}
% preamble.tex — shared setup for human-verification documents
% Compile each document with:  tectonic <File>.tex

\usepackage[T1]{fontenc}
\usepackage[utf8]{inputenc}
\usepackage{lmodern}
\usepackage{amssymb,amsmath,amsthm}
\usepackage{mathtools}
\usepackage[dvipsnames]{xcolor}
\usepackage{listings}
\usepackage{hyperref}
\usepackage[margin=2.5cm]{geometry}
\usepackage{parskip}

% ── Colours ──────────────────────────────────────────────────────────
\definecolor{leanBg}{HTML}{F7F7F0}
\definecolor{haskellBg}{HTML}{F0F4F7}
\definecolor{keyword}{HTML}{7B3294}
\definecolor{comment}{HTML}{5A8C32}
\definecolor{noteFg}{HTML}{B35900}

% ── Lean 4 listings language ─────────────────────────────────────────
% NOTE: Lean listing blocks use ASCII approximations of unicode symbols
% (e.g. -> for arrow, o. for compose).  The literate key substitutes
% these with proper math-mode renderings in the PDF.
\lstdefinelanguage{Lean4}{
  morekeywords={def,theorem,structure,where,namespace,end,import,fun,
                let,have,show,by,exact,intro,cases,rfl,subst,obtain,
                funext,simp,rw,apply,variable,Type,Prop,instance,
                section,open,class,extends,deriving,sorry,noncomputable},
  sensitive=true,
  morecomment=[l]{--},
  morecomment=[n]{/-}{-/},
  morestring=[b]",
}

% ── Lean listing style ───────────────────────────────────────────────
\lstdefinestyle{lean}{
  language=Lean4,
  backgroundcolor=\color{leanBg},
  basicstyle=\small\ttfamily,
  keywordstyle=\color{keyword}\bfseries,
  commentstyle=\color{comment}\itshape,
  breaklines=true,
  frame=single,
  framerule=0.4pt,
  xleftmargin=4pt,
  xrightmargin=4pt,
  aboveskip=6pt,
  belowskip=4pt,
  literate=
    {->}{{$\to$}}2
    {=>}{{$\Rightarrow$}}2
    {<|}{{$\langle$}}1
    {|>}{{$\rangle$}}1
    {==>}{{$\Longrightarrow$}}3
    { o. }{{$\;\circ\;$}}3
    {forall}{{$\forall$}}1
    {:=}{{$\,:=\,$}}2
    {!=}{{$\neq$}}2
    {<=}{{$\leq$}}2
    {>=}{{$\geq$}}2,
}

% ── Haskell listing style ───────────────────────────────────────────
\lstdefinestyle{haskell}{
  language=Haskell,
  backgroundcolor=\color{haskellBg},
  basicstyle=\small\ttfamily,
  keywordstyle=\color{keyword}\bfseries,
  commentstyle=\color{comment}\itshape,
  breaklines=true,
  frame=single,
  framerule=0.4pt,
  xleftmargin=4pt,
  xrightmargin=4pt,
  aboveskip=6pt,
  belowskip=4pt,
  literate=
    {->}{{$\to$}}2
    {=>}{{$\Rightarrow$}}2,
}

% ── Theorem environments ────────────────────────────────────────────
\theoremstyle{definition}
\newtheorem{defn}{Definition}[section]
\newtheorem{thm}[defn]{Theorem}
\newtheorem{prop}[defn]{Proposition}

\theoremstyle{remark}
\newtheorem*{fidelitynote}{Fidelity note}

% ── Convenience macros ───────────────────────────────────────────────
\newcommand{\lean}[1]{\texttt{#1}}      % inline Lean identifier
\newcommand{\hask}[1]{\texttt{#1}}      % inline Haskell identifier
\newcommand{\Ext}{\mathrm{Ext}}
\newcommand{\Pos}{\mathrm{Pos}}
\newcommand{\Shape}{\mathrm{Shape}}


\title{Human Verification: \texttt{Containers.Aggregation}\\[4pt]
  \large Aggregation as a Routing Cofunctor}
\author{Machine-generated from \texttt{lean/Aggregation.lean} for review by Neil Ghani}
\date{July 2026}

\begin{document}
\maketitle

\begin{abstract}
This document presents every definition and theorem \emph{statement} from the
Lean~4 module \lean{Containers.Aggregation} in three parallel forms: the
original Lean, a readable Haskell rendering, and precise mathematical English.
Proofs are omitted --- they are machine-checked by Lean.  The purpose is to
let a human verifier confirm that the formalisation captures the intended
mathematics.

The file verifies Construction~\texttt{con:aggregation} (\S9 of the paper): a
\textbf{routing policy} is a cofunctor (directed-container morphism) from the
concrete 11-facility Sawasdee Cellars network to a 2-node aggregate network
(Australia $\to$ Thailand).
\end{abstract}

\tableofcontents
\bigskip

%===================================================================
\section{Overview}
%===================================================================

A \emph{directed-container morphism} (cofunctor) from $\mathcal{D}$ to
$\mathcal{E}$ consists of:
\begin{itemize}
  \item a map on shapes $u : \mathcal{D}.\mathrm{Shape} \to \mathcal{E}.\mathrm{Shape}$, and
  \item a \emph{contravariant} position lift:
    $\ell : (s : \mathcal{D}.\mathrm{Shape}) \to \mathcal{E}.\mathrm{Pos}(u(s)) \to \mathcal{D}.\mathrm{Pos}(s)$,
\end{itemize}
satisfying three laws M0, M1, M2 (shape compatibility, root preservation, shift
compatibility).

Here $\mathcal{D}$ is the Sawasdee Cellars network and $\mathcal{E}$ is the
2-node aggregate network.  The shape map sends each facility to its region
(Australia or Thailand); the lift sends each aggregate route to a canonical
concrete route for that facility.


%===================================================================
\section{The aggregate network}
%===================================================================

% --- Region --------------------------------------------------------
\subsection{Inductive type: \lean{Region}}

\begin{defn}[Aggregate regions]
The 2-node aggregate graph has two nodes: \lean{australia} and \lean{thailand},
representing the two geographical regions of the wine supply chain.
\end{defn}

\paragraph{Lean.}
\begin{lstlisting}[style=lean]
inductive Region | australia | thailand
  deriving DecidableEq, Repr
\end{lstlisting}

\paragraph{Haskell.}
\begin{lstlisting}[style=haskell]
data Region = Australia | Thailand deriving (Eq, Show)
\end{lstlisting}


% --- AgLeg ---------------------------------------------------------
\subsection{Inductive type: \lean{AgLeg}}

\begin{defn}[Aggregate leg]
There is a single aggregate leg \lean{au2th} representing the sea-freight
connection from Australia to Thailand.
\end{defn}

\paragraph{Lean.}
\begin{lstlisting}[style=lean]
inductive AgLeg | au2th
  deriving DecidableEq, Repr
\end{lstlisting}

\paragraph{Haskell.}
\begin{lstlisting}[style=haskell]
data AgLeg = Au2th deriving (Eq, Show)
\end{lstlisting}


% --- GAgg ----------------------------------------------------------
\subsection{Definition: \lean{GAgg}}

\begin{defn}[2-node aggregate digraph]
$G_{\mathrm{agg}}$ is the digraph with $V = \{\mathrm{australia},
\mathrm{thailand}\}$, $E = \{\mathrm{au2th}\}$, $\mathrm{src}(\mathrm{au2th})
= \mathrm{australia}$, $\mathrm{tgt}(\mathrm{au2th}) = \mathrm{thailand}$.

Its network container $\mathcal{N}(G_{\mathrm{agg}})$ has:
\begin{itemize}
  \item $\mathrm{Pos}(\mathrm{australia}) = \{(\mathrm{australia}, \varepsilon),\;
    (\mathrm{thailand}, \mathrm{au2th})\}$ --- two aggregate routes;
  \item $\mathrm{Pos}(\mathrm{thailand}) = \{(\mathrm{thailand}, \varepsilon)\}$
    --- one aggregate route (sink).
\end{itemize}
\end{defn}

\paragraph{Lean.}
\begin{lstlisting}[style=lean]
@[reducible] def GAgg : Digraph where
  V := Region
  E := AgLeg
  src := agSrc
  tgt := agTgt
\end{lstlisting}

\paragraph{Haskell.}
\begin{lstlisting}[style=haskell]
gAgg :: Digraph Region AgLeg
gAgg = Digraph { src = \Au2th -> Australia
               , tgt = \Au2th -> Thailand }
\end{lstlisting}


%===================================================================
\section{The routing policy}
%===================================================================

% --- region --------------------------------------------------------
\subsection{Definition: \lean{region} (shape map)}

\begin{defn}[Region assignment]
Each facility belongs to a region:
\begin{align*}
  \mathrm{region}(v) &= \mathrm{australia} \quad
    \text{for } v \in \{\mathrm{vineyard}_i,\, \mathrm{sawasdee},\, \mathrm{melbournePort}\}, \\
  \mathrm{region}(v) &= \mathrm{thailand} \quad
    \text{for } v \in \{\mathrm{thaiPort},\, \mathrm{distributor}_i,\, \mathrm{retail}_i\}.
\end{align*}
This is the shape map $u$ of the routing-policy cofunctor.
\end{defn}

\paragraph{Lean.}
\begin{lstlisting}[style=lean]
@[reducible] def region : Facility -> Region
  | .vineyard1 | .vineyard2 | .vineyard3 => australia
  | .sawasdee | .melbournePort           => australia
  | .thaiPort                            => thailand
  | .distributor1 | .distributor2        => thailand
  | .retail1 | .retail2 | .retail3       => thailand
\end{lstlisting}

\paragraph{Haskell.}
\begin{lstlisting}[style=haskell]
region :: Facility -> Region
region Vineyard1    = Australia
region Vineyard2    = Australia
region Vineyard3    = Australia
region Sawasdee     = Australia
region MelbournePort = Australia
region _            = Thailand
\end{lstlisting}


% --- lift ----------------------------------------------------------
\subsection{Definition: \lean{lift} (position lift)}

\begin{defn}[Canonical concrete route lift]
The \textbf{lift} $\ell(f, \rho)$ maps an aggregate route $\rho :
\mathcal{N}(G_{\mathrm{agg}}).\mathrm{Pos}(\mathrm{region}(f))$ to a
canonical concrete route out of facility $f$.

\begin{itemize}
  \item If $\rho = (\mathrm{australia}, \varepsilon)$ (stay in Australia):
    $\ell(f, \rho) = (f, \varepsilon)$ for any Australian facility $f$.
  \item If $\rho = (\mathrm{thailand}, \mathrm{au2th})$ (cross to Thailand):
    $\ell(f, \rho)$ is the canonical route from $f$ to Thai Port via the
    fixed outbound path for $f$.
  \item For Thai facilities, $\mathrm{Pos}(\mathrm{region}(f)) =
    \{(\mathrm{thailand}, \varepsilon)\}$, so $\ell(f, \rho) = (f, \varepsilon)$.
\end{itemize}
\end{defn}

\paragraph{Lean.}
\begin{lstlisting}[style=lean]
def lift : (f : Facility) -> AgPos (region f) -> Route f
  | .vineyard1, ar => match ar with
      | <|australia, .nil _|> => <|_, .nil _|>
      | <|thailand,  _|>      =>
          alongLeg Leg.v1s (alongLeg Leg.smp (alongLeg Leg.sea <|_, .nil _|>))
  -- ... (similarly for all 11 facilities)
\end{lstlisting}

\paragraph{Haskell.}
\begin{lstlisting}[style=haskell]
lift :: Facility -> AgRoute Region AgLeg -> Route Facility Leg
lift Vineyard1 (Australia, []) = (Vineyard1, [])
lift Vineyard1 (Thailand, _)   = (ThaiPort, [V1s, Smp, Sea])
lift Vineyard2 (Australia, []) = (Vineyard2, [])
lift Vineyard2 (Thailand, _)   = (ThaiPort, [V2s, Smp, Sea])
-- ... similarly for all facilities
lift f         (Thailand, [])  = (f, [])  -- Thai facilities: stay put
\end{lstlisting}

\begin{fidelitynote}
The Lean type of \lean{lift} is genuinely dependent: the aggregate position
type is $\mathcal{N}(G_{\mathrm{agg}}).\mathrm{Pos}(\mathrm{region}(f))$,
which varies with $f$.  Lean can compute $\mathrm{region}(f)$ concretely by
case-splitting on $f$, which is why the definition pattern-matches on $f$
before matching on the aggregate route.
\end{fidelitynote}


% --- Law M0 --------------------------------------------------------
\subsection{Theorem: \lean{routingM0} (Law M0 --- shape compatibility)}

\begin{thm}[M0: shape compatibility]
For all facilities $f$ and aggregate routes $\rho :
\mathcal{N}(G_{\mathrm{agg}}).\mathrm{Pos}(\mathrm{region}(f))$:
\[
  \mathrm{region}\bigl(\mathcal{N}(\mathrm{graph}).\mathrm{sub}(f,\; \ell(f, \rho))\bigr)
  \;=\;
  \mathcal{N}(G_{\mathrm{agg}}).\mathrm{sub}\bigl(\mathrm{region}(f),\; \rho\bigr).
\]
The endpoint of the concrete route (in the wine network) maps to the
aggregate endpoint under \lean{region}.  Proved by cases on $f$ and $\rho$.
\end{thm}

\paragraph{Lean.}
\begin{lstlisting}[style=lean]
theorem routingM0 (f : Facility) (ar : AgPos (region f)) :
    region (Wine.network.sub f (lift f ar)) =
    (Digraph.network GAgg).sub (region f) ar
\end{lstlisting}

\paragraph{Haskell.}
\begin{lstlisting}[style=haskell]
-- M0: region (sub_wine f (lift f ar)) = sub_agg (region f) ar
-- For all f, ar: the endpoint of the concrete route has the right region.
propM0 :: Facility -> AgRoute Region AgLeg -> Bool
propM0 f ar = region (endpoint (lift f ar)) == agEndpoint (region f) ar
\end{lstlisting}


% --- Law M1 --------------------------------------------------------
\subsection{Theorem: \lean{routingM1} (Law M1 --- root preservation)}

\begin{thm}[M1: root preservation]
For all facilities $f$:
\[
  \ell\!\left(f,\; \mathcal{N}(G_{\mathrm{agg}}).\mathrm{root}(\mathrm{region}(f))\right)
  \;=\;
  \mathcal{N}(\mathrm{graph}).\mathrm{root}(f).
\]
Lifting the aggregate empty route gives the concrete empty route.
\end{thm}

\paragraph{Lean.}
\begin{lstlisting}[style=lean]
theorem routingM1 (f : Facility) :
    lift f ((Digraph.network GAgg).root (region f)) =
    Wine.network.root f
\end{lstlisting}

\paragraph{Haskell.}
\begin{lstlisting}[style=haskell]
-- M1: lift f (agRoot (region f)) = wineRoot f
-- Lifting the aggregate empty route is the concrete empty route.
propM1 :: Facility -> Bool
propM1 f = lift f (agRoot (region f)) == wineRoot f
\end{lstlisting}


% --- Law M2 --------------------------------------------------------
\subsection{Theorem: \lean{routingM2} (Law M2 --- shift compatibility)}

\begin{thm}[M2: shift compatibility]
For all $f : \mathrm{Facility}$, aggregate routes $\rho :
\mathcal{N}(G_{\mathrm{agg}}).\mathrm{Pos}(\mathrm{region}(f))$, and
$k : \mathcal{N}(G_{\mathrm{agg}}).\mathrm{Pos}(\mathcal{N}(G_{\mathrm{agg}}).\mathrm{sub}(\mathrm{region}(f), \rho))$:
\[
  \ell\bigl(f,\; \mathcal{N}(G_{\mathrm{agg}}).\mathrm{shift}(\mathrm{region}(f), \rho, k)\bigr)
  \;=\;
  \mathcal{N}(\mathrm{graph}).\mathrm{shift}\!\left(f,\; \ell(f, \rho),\;
    \ell\!\left(\mathcal{N}(\mathrm{graph}).\mathrm{sub}(f, \ell(f,\rho)),\;
    \mathrm{M0}^{-1}(k)\right)\right).
\]
Lifting a concatenated aggregate route equals the concrete concatenation of
the lifted segments.
\end{thm}

\paragraph{Lean.}
\begin{lstlisting}[style=lean]
theorem routingM2 (f : Facility) (ar : AgPos (region f))
    (k : (Digraph.network GAgg).Pos
           ((Digraph.network GAgg).sub (region f) ar)) :
    lift f ((Digraph.network GAgg).shift (region f) ar k) =
    Wine.network.shift f (lift f ar)
      (lift (Wine.network.sub f (lift f ar))
            ((routingM0 f ar).symm cast k))
\end{lstlisting}

\begin{fidelitynote}
The transport \lean{(routingM0 f ar).symm cast k} (Lean notation: \texttt{.symm \(\blacktriangleright\) k}) coerces $k$ along the M0
equality.  In the concrete proof, M0 is \lean{rfl} in all cases, so the
transport reduces to the identity.  The Haskell rendering omits this coercion
entirely.
\end{fidelitynote}

\paragraph{Haskell.}
\begin{lstlisting}[style=haskell]
-- M2: lift f (agShift (region f) ar k)
--     = wineShift f (lift f ar) (lift (wineSub f (lift f ar)) k)
-- (Omitting the transport; it is trivial in all concrete cases.)
propM2 :: Facility -> AgRoute Region AgLeg
       -> AgRoute Region AgLeg -> Bool
propM2 f ar k =
  lift f (agShift (region f) ar k) ==
  wineShift f (lift f ar) (lift (wineSub f (lift f ar)) k)
\end{lstlisting}


% --- routingPolicy -------------------------------------------------
\subsection{Definition: \lean{routingPolicy}}

\begin{defn}[Routing policy as a directed-container morphism]
The \textbf{routing policy} is the directed-container morphism (cofunctor)
\[
  \mathrm{routingPolicy} : \mathrm{DContMorphism}(\mathcal{N}(\mathrm{graph}),\;
  \mathcal{N}(G_{\mathrm{agg}}))
\]
with shape map \lean{region} and position lift \lean{lift}, satisfying laws M0,
M1, M2.  This verifies Construction~\texttt{con:aggregation}.
\end{defn}

\paragraph{Lean.}
\begin{lstlisting}[style=lean]
def routingPolicy : DContMorphism Wine.network (Digraph.network GAgg) where
  onShapes := region
  onPos    := lift
  m0       := routingM0
  m1       := routingM1
  m2       := routingM2
\end{lstlisting}

\paragraph{Haskell.}
\begin{lstlisting}[style=haskell]
data DContMorphism d e = DContMorphism
  { onShapes :: Shape d -> Shape e
  , onPos    :: (s :: Shape d) -> Pos e (onShapes s) -> Pos d s
  -- plus proofs of M0, M1, M2
  }

routingPolicy :: DContMorphism WineNetwork AggNetwork
routingPolicy = DContMorphism
  { onShapes = region
  , onPos    = lift
  }
\end{lstlisting}

\begin{fidelitynote}
A directed-container morphism is a \emph{cofunctor}: the position map goes
\emph{backwards}, from aggregate positions to concrete positions, just as a
container morphism's position map is contravariant.  The shape map goes
forward (facility $\to$ region) while the position lift goes backward (aggregate
route $\to$ concrete route).
\end{fidelitynote}


%===================================================================
\section{Translation conventions}
%===================================================================

\begin{enumerate}
\item \textbf{Dependent position types.}  The type of \lean{lift} is
  $\Pi(f : \mathrm{Facility}),\, \mathcal{N}(G_{\mathrm{agg}}).\mathrm{Pos}(\mathrm{region}(f)) \to \mathrm{Route}(f)$.
  The first argument determines the type of the aggregate position; this is a
  genuinely dependent type that Haskell cannot express directly.

\item \textbf{Transport in M2.}  The cast \lean{(routingM0 f ar).symm cast k} (Lean: \texttt{.symm \(\blacktriangleright\) k})
  appears because the endpoint of \lean{lift f ar} has type \lean{Facility},
  while $k$'s domain is parametrised by a \lean{Region}.  The M0 equality
  converts between them.  In all concrete cases, M0 is \lean{rfl} and the
  transport is the identity.

\item \textbf{Proof bodies omitted.}  Only definitions and theorem
  \emph{statements} are translated.  Lean checks the proofs; the human
  verifier checks that the \emph{statements} match the intended mathematics.
\end{enumerate}

\end{document}
